\section*{Lectures on Lie Groups and Lie Algebras – 6,7}

\section*{The Lie Algebra $\mathfrak{sl}_2$}

Let $e,f,h$ be the standard generators of $\mathfrak{sl}_2(\mathbb{C})$ such that
\[
[e,f] = h, \quad [h,e] = 2e, \quad [h,f] = -2f.
\]

\textbf{Proposition 1.}
\[
\mathfrak{sl}_2(\mathbb{C}) \cong \mathfrak{so}_3(\mathbb{C}).
\]

\textbf{Proof.}
On $\mathfrak{sl}_2(\mathbb{C}) \cong \mathbb{C}^3$ there exists a nondegenerate invariant symmetric bilinear form (the Killing form). Thus the image of $\mathfrak{sl}_2(\mathbb{C})$ under the adjoint representation lies in $\mathfrak{so}_3(\mathbb{C})$. Since the adjoint representation is faithful and $\mathfrak{so}_3(\mathbb{C})$ is 3-dimensional, $\mathrm{ad}$ is an isomorphism.
\qed

A natural basis of $\mathfrak{so}_3(\mathbb{C})$ is

\[
i =
\begin{pmatrix}
0&1&0\\
-1&0&0\\
0&0&0
\end{pmatrix}, \quad
j =
\begin{pmatrix}
0&0&0\\
0&0&1\\
0&-1&0
\end{pmatrix}, \quad
k =
\begin{pmatrix}
0&0&1\\
0&0&0\\
-1&0&0
\end{pmatrix}.
\]

with commutation relations
\[
[i,j]=k, \quad [j,k]=i, \quad [k,i]=j.
\]

Under the isomorphism $\mathfrak{sl}_2(\mathbb{C}) \cong \mathfrak{so}_3(\mathbb{C})$
these correspond to the Pauli matrices:

\[
i=
\begin{pmatrix}
0&-1\\
1&0
\end{pmatrix}, \quad
j=
\begin{pmatrix}
i&0\\
0&-i
\end{pmatrix}, \quad
k=
\begin{pmatrix}
0&i\\
i&0
\end{pmatrix}.
\]

\textbf{Proposition 2.}
\[
\mathfrak{sl}_2(\mathbb{R}) \not\cong \mathfrak{so}_3(\mathbb{R}).
\]

\textbf{Proof.}
Their Killing forms have different signatures:
$\mathfrak{sl}_2(\mathbb{R})$ has signature $(++-)$,
while $\mathfrak{so}_3(\mathbb{R})$ is negative definite.
\qed

\section*{Finite-Dimensional Irreducible Representations of $\mathfrak{sl}_2$}

Recall that
\[
S^n(\mathbb{C}^2)
\]
is irreducible.

\textbf{Theorem 1.}
Every finite-dimensional irreducible representation of
$\mathfrak{sl}_2(\mathbb{C})$ is isomorphic to $S^n(\mathbb{C}^2)$
for some $n\in\mathbb{Z}_{\ge0}$.

\subsection*{Key Lemmas}

If $v$ is an eigenvector of $\rho(h)$ with eigenvalue $\mu$, then
\[
\rho(e)v \text{ has eigenvalue } \mu+2,
\quad
\rho(f)v \text{ has eigenvalue } \mu-2.
\]

There exists a highest-weight vector $v_\lambda$ such that
\[
\rho(e)v_\lambda=0.
\]

The space
\[
V_\lambda = \mathrm{span}\{ \rho(f)^k v_\lambda \}
\]
is a subrepresentation.

Finite dimensionality forces
\[
\lambda = n\in \mathbb{Z}_{\ge0}.
\]

Thus irreducible representations are exactly $V_n$,
and $\dim V_n=n+1$.

\section*{Category $\mathcal{O}$}

Finite-dimensional irreducible representations satisfy:

1. $\rho(h)$ is semisimple.  
2. $\rho(e)$ acts locally nilpotently.

The category of finitely generated $\mathfrak{sl}_2$-modules satisfying these properties
is called category $\mathcal{O}$.

\textbf{Verma module $M_\lambda$:}

Basis $v_\lambda^{(k)}, k=0,1,2,\dots$ with

\[
\rho(e)v_\lambda^{(k)} = (k\lambda-k(k-1))v_\lambda^{(k-1)},
\]
\[
\rho(f)v_\lambda^{(k)} = v_\lambda^{(k+1)},
\]
\[
\rho(h)v_\lambda^{(k)} = (\lambda-2k)v_\lambda^{(k)}.
\]

It is reducible iff $\lambda\in\mathbb{Z}_{\ge0}$.

\section*{Casimir Operator and Complete Reducibility}

Define
\[
C = ef + fe + \frac12 h^2.
\]

\textbf{Lemma.}
$C$ lies in the center of $U(\mathfrak{sl}_2)$.

By Schur’s lemma, $C$ acts by scalars on irreducibles.

\[
\rho(C)|_{V_n} = \frac{n(n+2)}{2}\,\mathrm{id}.
\]

\textbf{Theorem 2.}
Every finite-dimensional representation of
$\mathfrak{sl}_2(\mathbb{C})$ is completely reducible.

\section*{Characters}

Let $V(m)$ be the weight space of weight $m$.

\[
\chi_V(q)=\sum_{m\in\mathbb{Z}} q^m \dim V(m).
\]

Properties:
\[
\chi_{V\oplus W}=\chi_V+\chi_W,
\quad
\chi_{V\otimes W}=\chi_V\chi_W.
\]

\[
\chi_{V_n}(q)
=
q^n+q^{n-2}+\cdots+q^{-n}
=
\frac{q^{n+1}-q^{-n-1}}{q-q^{-1}}.
\]

Example:
\[
\chi_{V_2\otimes V_2}
=
(q^2+1+q^{-2})^2
=
\chi_{V_4}+\chi_{V_2}+\chi_{V_0}.
\]

Thus
\[
V_2\otimes V_2 \cong V_4\oplus V_2\oplus V_0.
\]